\section*{Chapitre \ref{ch:b3:hydrogen-atom} --- L'atome d'hydrogène}

\begin{solution}{exo:b3:hydrogen-atom:1}
(a) $u'' = (1/a^2 - 2/ar)u$~; les termes en $1/r$ s'identifient quand
$\hbar^2/m_{\text{e}}a = k$, c'est-à-dire $a = a_0$, et les termes constants
donnent $E = -\hbar^2/2m_{\text{e}}a_0^2 = -E_{\text{I}}$. (b) $C =
2/a_0^{3/2}$. (c) $[\hbar^2/mk] = \unit{m}$~; $[mk^2/\hbar^2] =
\unit{J}$. (d) L'argument variationnel de
\cref{exo:b3:schrodinger-three-dimensions:12} montre que
$-E_{\text{I}}$ est la plus petite énergie qu'un état normé puisse
atteindre~: le principe d'indétermination met un plancher à l'atome.
\end{solution}

\begin{solution}{exo:b3:hydrogen-atom:2}
(a) \qty{121.6}{nm}~; limite \qty{91.2}{nm}. (b) $656.3$ (rouge),
$486.1$ (bleu-vert), \qty{434.0}{nm} (violet)~; limite
\qty{364.6}{nm} --- les trois premières sont visibles. (c) De
\qty{1875}{nm} jusqu'à \qty{820}{nm}~: proche infrarouge. (d)
$R_{\text{H}}c\,(1/109^2 - 1/110^2) = \qty{5.01}{GHz}$.
\end{solution}

\begin{solution}{exo:b3:hydrogen-atom:3}
(a) $\dd(r^2\eu^{-2r/a_0})/\dd r = 0$ en $r = a_0$. (b) $\langle r\rangle = (4/a_0^3)\int_0^\infty
r^3\eu^{-2r/a_0}\dd r = (4/a_0^3)\times 3!\,(a_0/2)^4 = \tfrac32
a_0$. (c) Avec $x = 4$~:
$\tfrac12(16 + 8 + 2)\eu^{-4} = 0.24$~: un quart du temps, l'électron
est au-delà de $2a_0$. (d) L'électron n'a pas de rayon~: $P(r)$ est
une distribution avec un mode, une moyenne et de longues queues ---
l'atome est flou au facteur deux près.
\end{solution}

\begin{solution}{exo:b3:hydrogen-atom:4}
(a) $l = 0$~: un~; $l = 1$~: trois~; $l = 2$~: cinq --- neuf. (b) $2$,
$8$, $18$. (c) Exactement les longueurs des trois premières lignes~:
le tableau périodique est le registre de remplissage de ces couches.
(d) La dégénérescence en $m$ survit (l'espace reste isotrope)~; celle
en $l$ se brise parce que le potentiel écranté que voit l'électron de
valence n'est plus en $1/r$ pur ---
\cref{exo:b3:hydrogen-atom:8}.
\end{solution}

\begin{solution}{exo:b3:hydrogen-atom:5}
(a) Hydrogénoïde avec $Z - 1$~: $h\nu = (Z-1)^2E_{\text{I}}(1 -
\tfrac14)$. (b) $784 \times \qty{10.2}{eV} = \qty{8.0}{keV}$ ---
en accord avec la K$\alpha$ mesurée du cuivre. (c) Que l'entier qui
ordonne les éléments est la \emph{charge nucléaire}, et non la masse
atomique~: les lacunes dans les droites de Moseley ont localisé des
éléments non découverts. (d) $42^2 \times 10.2 = \qty{18.0}{keV}$ ---
le technétium, trouvé des décennies plus tard, s'y est conformé.
\end{solution}

\begin{solution}{exo:b3:hydrogen-atom:6}
(a) Deux masses égales~: $\mu = m_{\text{e}}/2$~: liaison
\qty{6.8}{eV}, rayon $2a_0 = \qty{0.106}{nm}$. (b) $\tfrac34
\times 6.8 = \qty{5.1}{eV}$~: \qty{243}{nm}. (c) Deux photons de
\qty{511}{keV}, dos à dos (conservation de la quantité de mouvement
au repos). (d) La tomographie par émission de positons~: la paire de
photons colinéaires de \qty{511}{keV} est le signal que triangule
chaque anneau TEP.
\end{solution}

\begin{solution}{exo:b3:hydrogen-atom:7}
(a) $a_{\text{ex}} = 0.0529 \times 12.9/0.058 = \qty{11.8}{nm}$.
(b) $E = 13.6 \times 0.058/12.9^2 = \qty{4.7}{meV}$. (c) La taille
naturelle de la paire dépasse celle de la boîte~: c'est la paroi, et
non l'attraction, qui façonne l'état --- le $1/R^2$ bat le $1/R$,
comme promis. (d) $k_{\text{B}}T$ à température ambiante vaut cinq
fois la liaison~: les paires s'ionisent~; la physique excitonique vit
au-dessous de $\sim\qty{50}{K}$ dans des cristaux propres.
\end{solution}

\begin{solution}{exo:b3:hydrogen-atom:8}
(a) $3s$~: son lobe le plus interne plonge dans le cœur, où le $+11$
nucléaire est à peine écranté. (b) $-13.6/9 = \qty{-1.51}{eV}$
correspond presque exactement à $3d$~: $3d$ orbite hors du cœur et ne
voit qu'un $+1$ net~; $3s$ (et en partie $3p$) goûtent le potentiel
profond et s'enfoncent. (c) $\Delta E = \qty{2.10}{eV}$~:
$\lambda = \qty{590}{nm}$ --- le jaune sodium des lampadaires. (d) $n_{\text{eff}} =
\sqrt{13.6/|E|}$~: $1.63$ et $2.11$, d'où $\delta_s = 1.37$,
$\delta_p = 0.89$.
\end{solution}

\begin{solution}{exo:b3:hydrogen-atom:9}
(a) À peu près la liaison de $n = 100$, $\sim\qty{1.4}{meV}$, plus la
faible énergie thermique de l'électron~: un photon d'infrarouge
lointain ou radio. (b) Les pas $\Delta n = 1$ vers $n \approx 100$
tombent dans la bande radio du GHz --- les raies de recombinaison. (c)
H$\alpha$ est l'étape $3 \to
2$, près de la fin de la cascade. (d) Lyman $\alpha$, quoique la plus
intense, est ultraviolette et, dans une nébuleuse pleine d'hydrogène
dans son état fondamental, elle est diffusée de façon résonnante et
piégée~; H$\alpha$ s'échappe librement et peint la nébuleuse en rouge.
\end{solution}

\begin{solution}{exo:b3:hydrogen-atom:10}
(a) $\langle 1/r\rangle = 1/a_0$~: $\langle E_p\rangle = -k/a_0 =
\qty{-27.2}{eV} = 2E_1$. (b) $\langle E_k\rangle = E_1 - \langle
E_p\rangle = +\qty{13.6}{eV}$~: le rapport du viriel $-\tfrac12$ de
toute orbite en $1/r$. (c) $v = \sqrt{2E_k/m_{\text{e}}} = \alpha c =
\qty{2.2e6}{m/s}$. (d) $I = ev/2\pi a_0 \approx \qty{1}{mA}$
circulant à \qty{53}{pm}~: $B \sim \mu_0I/2a_0 \approx
\qty{12}{T}$ --- l'énorme champ interne dont se nourrira la structure
hyperfine.
\end{solution}

\begin{solution}{exo:b3:hydrogen-atom:11}
(a) D'après \cref{exo:b3:hydrogen-atom:10}(c), $v/c = \alpha$~;
par réétalonnage~: $Z\alpha/n$. (b) $\alpha^2E_{\text{I}} = \qty{0.72}{meV}$~;
à $n = 2$ les écarts sortent à quelques dizaines de
$\unit{\micro eV}$ --- quelque \qty{10}{GHz}, mesurables par radio (et
mesurés). (c) $Z\alpha =
0.67$~: les deux tiers de la vitesse de la lumière --- le traitement
de Schrödinger échoue~; l'équation de Dirac prend le relais. (d) Quand
$Z\alpha \to 1$, l'énergie de l'état fondamental plonge formellement
sous $-2m_{\text{e}}c^2$~: le vide lui-même produirait des paires
électron--positon --- les champs supercritiques, cherchés dans les
collisions d'ions lourds.
\end{solution}

\begin{solution}{exo:b3:hydrogen-atom:12}
(a) Dériver~: $\dot{\vect p}\wedge\vect L = -(mk/r^3)\vect
r\wedge(\vect r\wedge\dot{\vect r}\,m)$~; en développant le double
produit vectoriel on obtient exactement $mk\,\dd\vect e_r/\dd t$~: $\dot{\vect A} =
\vect 0$. (b) Un axe conservé est une symétrie au-delà de
l'isotropie~; une dégénérescence est un étiquetage incomplet, et
l'étiquette supplémentaire (le cousin quantique de $\vect
A$) relie les $l$ d'un même $n$. (c) Avec
$V = -k/r + \epsilon/r^2$ le moment cinétique effectif se décale, la
période angulaire ne s'accorde plus à la période radiale, et l'axe de
l'ellipse tourne à un rythme $\propto\epsilon$. (d) Les corrections de
la relativité jouent le rôle d'$\epsilon$ dans l'hydrogène même~: la
structure fine de \cref{exo:b3:hydrogen-atom:11} est précisément la
brisure de la dégénérescence en $l$.
\end{solution}

\begin{solution}{pb:b3:hydrogen-atom:1}
\textbf{1.} $\langle r\rangle \sim n^2a_0$~; $|E_n| =
E_{\text{I}}/n^2$~; écart $\approx 2E_{\text{I}}/n^3$~; fréquence
orbitale $\propto 1/n^3$ (Kepler sur l'orbite de Bohr).
\textbf{2.} $E_{n+1} - E_n \to 2E_{\text{I}}/n^3$, et $h$ fois la
fréquence classique $v_n/2\pi r_n$ donne la même expression (le calcul
de \cref{pb:b3:hamiltonian-mechanics:1}, question 20).
\textbf{3.} $n = 50$~: $\qty{132}{nm}$, \qty{5.4}{meV}~; $n = 100$~:
$\qty{0.53}{\micro m}$, \qty{1.4}{meV}.
\textbf{4.} $n^2ea_0 = \qty{2.1e-26}{C.m} \approx 6300\,\text{D}$~:
l'équivalent de trois mille molécules d'eau de dipôle sur un seul atome.
\textbf{5.} $F \sim \num{5.4e-3}/\num{1.32e-7} \approx
\qty{4e4}{V/m} = \qty{400}{V/cm}$~; le seuil exact d'ionisation
classique est environ huit fois plus faible ($\sim\qty{50}{V/cm}$)~:
dans les deux cas, un souffle de champ.
\textbf{6.} Avec des sections efficaces géométriques
$\sim\pi n^4a_0^2$, même un vide de laboratoire ultra-poussé fait
entrer un tel atome en collision en quelques microsecondes~; les
densités interstellaires ($\sim\num{e6}$ particules par
$\unit{m^3}$) lui laissent des jours --- l'espace est la meilleure
enceinte à vide.
\textbf{7.} $\nu = R_{\text{H}}c\,(1/109^2 - 1/110^2) =
\qty{5.01}{GHz}$.
\textbf{8.} Radio centimétrique~: un radiotélescope --- une grande
antenne et un récepteur silencieux.
\textbf{9.} $\lambda = \qty{6.0}{cm}$~; l'atome $n = 110$ mesure
$\sim\qty{0.6}{\micro m}$~: cent mille atomes par longueur d'onde ---
confortablement un régime d'«~antenne~».
\textbf{10.} $\Delta\nu/\nu \sim v/c = \num{4.3e-5}$~: environ
\qty{200}{kHz} sur \qty{5}{GHz} --- et la largeur mesurée rend la
température nébulaire.
\textbf{11.} La Galaxie interne ionisée~: des régions HII cachées
derrière une poussière qui éteint tout photon de Balmer --- les raies
de recombinaison radio ont cartographié la structure spirale que
l'astronomie optique ne pouvait voir.
\textbf{12.} Les micro-champs électriques des ions voisins
(élargissement Stark) brouillent les niveaux, et le rayonnement du
corps noir ou cosmique ainsi que les collisions ionisent ou mélangent
en $l$ ces états fragiles.
\textbf{13.} $C_6 \sim d^4/\Delta E \propto (n^2)^4/n^{-3} =
n^{11}$~: porter $n$ de $1$ à $70$ et l'interaction croît de vingt
ordres de grandeur.
\textbf{14.} À l'intérieur de $R_{\text{b}}$, l'état de paire est
décalé hors résonance, si bien que l'excitation d'un atome
\emph{conditionne} la réponse de son voisin~: exactement la logique
contrôlée dont a besoin une porte à deux qubits.
\textbf{15.} Des interactions de portée micrométrique permettent à
chaque atome de siéger dans sa propre pince adressable, très éloigné
selon les critères atomiques et pourtant fortement couplé --- une
portée d'interaction accordée à la résolution optique.
\textbf{16.} $\sim 200$ portes par durée de vie~: ce rapport borne la
fidélité atteignable, et c'est lui, non les microsecondes brutes, dont
vit un processeur.
\textbf{17.} À \qty{300}{K} le bain de photons est dense vers
$\sim\qty{0.1}{meV}$~: il brasse et ionise les niveaux de Rydberg en
moins d'une durée de vie radiative~; des écrans froids (\qty{4}{K})
affament ces transitions.
\textbf{18.} Les taux varient comme $d^2 \propto n^4$~: à $n =
50$, six millions de fois le couplage d'un atome ordinaire ---
et c'est pourquoi une cellule de vapeur d'atomes de Rydberg est
aujourd'hui un capteur étalonné de champs micro-ondes et térahertz.
\textbf{19.} De $92^2 \times \qty{13.6}{eV} = \qty{115}{keV}$
jusqu'à $13.6/10^6\,\unit{eV} = \qty{13.6}{\micro eV}$~: dix ordres de
grandeur à partir d'une seule formule.
\textbf{20.} Le classique se lève là où l'écart devient une fraction
évanescente de l'énergie, $n \gg 1$ --- la limite de correspondance
propre à l'atome.
\textbf{21.} C'est le seul atome calculable à partir des premiers
principes jusqu'à la précision de l'expérience~: tout désaccord serait
une découverte, aussi l'hydrogène ancre-t-il les constantes
fondamentales.
\textbf{22.} $\tau \approx 1.6\,\unit{ns} \times 50^3 =
\qty{0.2}{ms}$ --- généreux, mais la redistribution par le corps noir
(partie III) l'entame, une raison de plus pour les écrans froids.
\textbf{23.} La symétrie CPT --- atomes de matière et d'antimatière
doivent coïncider raie par raie~; l'hydrogène est l'ancre parce que
lui seul voit la théorie atteindre le quinzième chiffre à côté de
l'expérience.
\textbf{24.} Le dipôle en $n^2$ fait sentir à l'atome le champ d'un
seul photon micro-onde, tandis que sa longue durée de vie et sa
structure de niveaux lui permettent d'acquérir une phase mesurable
\emph{sans} absorber le photon~: de la détection quantique non
destructive.
\textbf{25.} Les lois d'échelle $n^2$, $n^{-2}$, $n^{-3}$, $n^{11}$
étirent l'unique atome résolu de \qty{52.9}{pm} à un demi-micromètre
et de \qty{121.6}{nm} à \qty{6}{cm} --- la même solution de
Schrödinger émettant depuis les régions HII de la Galaxie et cadençant
des portes à deux qubits dans des réseaux de pinces optiques.
\end{solution}
